\documentclass[11pt]{article}
\usepackage[utf8]{inputenc}
\usepackage[margin=1in]{geometry}
\usepackage{amsmath}
\usepackage{graphicx}
\usepackage{booktabs}
\usepackage{hyperref}
\usepackage[round]{natbib}

\title{Auditory Cortex Conveys Non-Topographic Sound Location Information to Primary Visual Cortex}
\author{Author A$^{1}$, Author B$^{1,2}$, Author C$^{2}$ \\
\small $^{1}$Department of Neuroscience, Erasmus MC, Rotterdam, Netherlands \\
\small $^{2}$Netherlands Institute for Neuroscience, Amsterdam, Netherlands}
\date{}

\begin{document}
\maketitle

\begin{abstract}
Auditory cortex (AC) sends direct projections to primary visual cortex (V1), but whether these inputs carry spatial auditory information and whether this information is topographically organized to match V1's retinotopic map is unknown. Here we use dual-color two-photon calcium imaging to simultaneously record AC axon boutons in V1 layer 1 (GCaMP6s) and postsynaptic V1 neurons in layer 2/3 (jRGECO1a) in head-fixed mice presented with spatially controlled auditory and visual stimuli from a custom 39-position speaker-LED array. We find that AC inputs to V1 carry robust information about sound source location in both azimuth and elevation, with population decoding significantly exceeding chance. However, the auditory spatial receptive fields of AC boutons are not retinotopically matched to the visual receptive fields of underlying V1 neurons---AC inputs are non-topographically organized. AC boutons are tuned to sounds spanning both ipsilateral and contralateral hemifields, covering a wider spatial range than V1's visual field representation. In contrast, secondary visual area lateral (V2L) inputs to V1 show retinotopic visual receptive fields but minimal auditory spatial coding. Furthermore, audiovisual modulation of V1 neuron responses does not depend on the spatial congruency of sound and light stimuli. These findings reveal that AC provides V1 with a distributed, non-topographic representation of auditory space that modulates visual processing independently of spatial alignment.
\end{abstract}

\section{Introduction}

The ability to bind spatially and temporally congruent stimuli from different sensory modalities is fundamental to coherent perception \citep{stein2012new}. Auditory cortex sends direct projections to primary visual cortex across species, and these projections are thought to mediate audiovisual interactions observed in V1 neurons \citep{iurilli2012sound, ibrahim2016cross}. In the superior colliculus, auditory and visual spatial maps are topographically aligned, providing a substrate for binding spatially congruent multisensory stimuli \citep{king2004superior}. Whether a similar topographic organization exists in cortex---specifically, whether AC inputs to V1 carry spatial information arranged to match V1's retinotopic map---has been a fundamental open question.

AC neurons show tuning to sound location, but it was unknown whether the subpopulation of AC neurons projecting to V1 specifically relays spatial auditory information \citep{stecker2005sound}. Previous work demonstrated that V1-projecting AC neurons selectively encode certain acoustic features, but spatial coding in these projections had not been characterized \citep{ibrahim2016cross}. If AC inputs to V1 were retinotopically arranged, this would provide a cortical mechanism for binding spatially congruent audiovisual stimuli analogous to the subcortical mechanism in the superior colliculus.

Here we address this question by combining a custom speaker-LED stimulus array with dual-color two-photon calcium imaging to simultaneously record presynaptic AC (or V2L) boutons and postsynaptic V1 neurons. This approach enables direct comparison of auditory spatial tuning in AC inputs with the visual receptive fields of their target neurons.

\section{Methods}

\subsection{Stimulus Array}

A custom speaker-LED array was constructed with 39 positions arranged in 13 columns (azimuth) and 3 rows (elevation), with 10-degree azimuth steps and 20-degree elevation steps. The array covered 20 degrees ipsilateral to 100 degrees contralateral in azimuth and $\pm$20 degrees in elevation relative to the mouse's head position. Auditory stimuli were white noise bursts (2--20 kHz, 1 second duration) and visual stimuli were flashing white LEDs (1 second duration). Table~\ref{tab:array} summarizes the stimulus parameters.

\begin{table}[htbp]
\centering
\caption{Stimulus array parameters.}
\label{tab:array}
\begin{tabular}{ll}
\toprule
Parameter & Value \\
\midrule
Total positions & 39 \\
Azimuth columns & 13 \\
Elevation rows & 3 \\
Azimuth step & 10 degrees \\
Elevation step & 20 degrees \\
Ipsilateral coverage & 20 degrees \\
Contralateral coverage & 100 degrees \\
Elevation range & $-$20 to $+$20 degrees \\
Auditory stimulus & White noise, 2--20 kHz, 1 s \\
Visual stimulus & Flashing white LED, 1 s \\
\bottomrule
\end{tabular}
\end{table}

\subsection{Animals and Imaging}

Two-photon calcium imaging was performed in head-fixed mice using a dual-color configuration. AC axon boutons in V1 layer 1 were labeled with GCaMP6s (green channel) via viral injection into AC ($n = 7$ mice, 42 sessions) or V2L ($n = 3$ mice, 15 sessions). Postsynaptic V1 neurons in layer 2/3 were labeled with jRGECO1a (red channel). This dual-color approach enabled simultaneous recording of presynaptic inputs and postsynaptic targets (Table~\ref{tab:animals}).

\begin{table}[htbp]
\centering
\caption{Experimental animals and imaging sessions.}
\label{tab:animals}
\begin{tabular}{lllll}
\toprule
Condition & $N$ mice & Sessions & Axon indicator & Soma indicator \\
\midrule
AC-injected & 7 & 42 & GCaMP6s & jRGECO1a \\
V2L-injected & 3 & 15 & GCaMP6s & jRGECO1a \\
\bottomrule
\end{tabular}
\end{table}

\subsection{Analysis}

A spatial modulation index (SMI) was computed for each bouton to quantify spatial selectivity. Population decoding used a maximum-likelihood framework to estimate sound location from the activity of AC bouton populations, with cross-validation and shuffled controls. Additional experiments with a single speaker spanning $-$90 to $+$90 degrees azimuth (2--80 kHz white noise) tested full azimuthal tuning. Audiovisual congruency experiments placed an LED within the visual receptive field of V1 neurons while varying speaker position to test whether V1 modulation depended on spatial match.

\section{Results}

\subsection{AC Inputs Are Predominantly Auditory, V2L Inputs Are Visual}

Comparison of response modalities between AC and V2L boutons in V1 revealed a striking dissociation. AC boutons showed a significantly higher fraction of auditory-responsive elements compared to visual-responsive elements (two-sided paired $t$-test, $p < 0.05$). In contrast, V2L boutons showed the opposite pattern, with visual responses dominating over auditory responses (Table~\ref{tab:modality}).

\begin{table}[htbp]
\centering
\caption{Fraction of responsive boutons by modality.}
\label{tab:modality}
\begin{tabular}{llll}
\toprule
Input Source & Auditory (\%) & Visual (\%) & Dominant Modality \\
\midrule
AC boutons & High (majority) & Low (minority) & Auditory \\
V2L boutons & Low (minority) & High (majority) & Visual \\
\bottomrule
\end{tabular}
\end{table}

\subsection{Population Decoding of Sound Location from AC Inputs}

Maximum-likelihood decoding analysis revealed that AC bouton populations in V1 carried significant information about sound source location. Decoding accuracy for both azimuth and elevation significantly exceeded chance levels (shuffled control) for AC inputs (two-sided paired $t$-test, $p < 0.05$). In contrast, V2L inputs showed near-chance decoding accuracy for both spatial dimensions (Table~\ref{tab:decoding}).

\begin{table}[htbp]
\centering
\caption{Sound location decoding accuracy comparison.}
\label{tab:decoding}
\begin{tabular}{llll}
\toprule
Dimension & AC Inputs & V2L Inputs & Shuffled Control \\
\midrule
Azimuth & Above chance$^{*}$ & Near chance & Baseline \\
Elevation & Above chance$^{*}$ & Near chance & Baseline \\
\bottomrule
\end{tabular}
\end{table}

\subsection{AC Inputs Lack Retinotopic Organization}

Despite carrying spatial auditory information, AC inputs to V1 were not retinotopically organized. The azimuth peaks of AC bouton auditory receptive fields did not correlate with the visual receptive field centers of underlying V1 neurons. AC auditory receptive fields spanned a wider range than V1's visual field coverage, including substantial representation of ipsilateral sounds well beyond the visual representation. In contrast, V2L inputs showed retinotopic visual receptive fields that matched the V1 retinotopic map (Table~\ref{tab:organization}).

\begin{table}[htbp]
\centering
\caption{Comparison of spatial organization properties.}
\label{tab:organization}
\begin{tabular}{llll}
\toprule
Property & AC Inputs & V1 Somata & V2L Inputs \\
\midrule
Spatial specificity & Broad & Sharp & Sharp \\
Total span & Larger than V1 & Contralateral & Contralateral \\
Ipsilateral coverage & Substantial & Minimal & Minimal \\
Topographic organization & Non-topographic & Retinotopic & Retinotopic \\
\bottomrule
\end{tabular}
\end{table}

\subsection{Broad Azimuthal Tuning Including Ipsilateral Space}

Extended azimuth experiments ($-$90 to $+$90 degrees, 2--80 kHz white noise) confirmed that individual AC boutons in V1 were tuned to sounds in both ipsilateral and contralateral hemifields. Some boutons preferred ipsilateral locations, others contralateral locations, and some showed broad tuning across the full azimuthal range. A substantial proportion of spatially selective AC boutons preferred ipsilateral sound positions, representing auditory space well outside V1's visual field representation.

\subsection{Audiovisual Modulation Is Independent of Spatial Congruency}

The audiovisual congruency experiment tested whether V1 neuron responses to LED flashes (positioned within each neuron's visual receptive field at 40--60 degrees azimuth) were differentially modulated by concurrent sounds presented at matched versus mismatched spatial positions ($-$40 to $+$40 degrees, 20-degree steps). V1 neuron responses were modulated by concurrent sounds, but this modulation did not depend on the spatial alignment between sound and light. The statistical analysis revealed no significant effect of spatial congruency on the magnitude of audiovisual modulation.

\subsection{Axon Distribution Confirms Layer 1 Targeting}

Histological analysis confirmed that both AC and V2L axons in V1 were densest in layer 1, with decreasing density in deeper layers. Normalized fluorescence intensity profiles from $n = 7$ AC-injected and $n = 3$ V2L-injected mice showed peaks in layer 1 for both projection sources, consistent with targeting of the apical dendritic tufts of pyramidal neurons.

\section{Discussion}

\subsection{Distributed Auditory Spatial Code in V1}

Our findings reveal that AC provides V1 with a distributed, non-topographic representation of auditory space. Unlike the topographically aligned auditory-visual maps in the superior colliculus \citep{king2004superior}, the cortical pathway from AC to V1 does not preserve spatial register between auditory and visual reference frames. Instead, individual AC boutons in V1 carry information about sound location across a broad spatial range, and the population as a whole can accurately estimate sound position through maximum-likelihood decoding.

This non-topographic organization has important implications for theories of multisensory binding. The classical view that binding requires spatial alignment \citep{stein2012new} may apply to subcortical structures but not to cortical cross-modal pathways. In cortex, the distributed spatial code may serve a different function, such as providing contextual information about the broader auditory environment rather than supporting precise audiovisual spatial binding.

\subsection{Functional Implications of Non-Congruency-Dependent Modulation}

The finding that audiovisual modulation of V1 neurons does not depend on spatial congruency is consistent with the non-topographic organization of AC inputs. If AC inputs were arranged to match V1 retinotopy, one would expect stronger modulation when sound and light originate from the same location. The absence of this effect suggests that AC-to-V1 projections modulate visual processing in a spatially global manner, potentially serving to enhance overall visual sensitivity or arousal in the presence of salient sounds regardless of their location \citep{kayser2007functional}.

\subsection{Contrast with V2L Inputs}

The complementary properties of AC and V2L inputs to V1---AC providing auditory spatial information without retinotopic organization, V2L providing retinotopic visual information without auditory spatial coding---suggest that V1 integrates qualitatively different types of cross-modal information from different cortical sources. This division of function may enable V1 to simultaneously benefit from spatially precise visual feedback (via V2L) and spatially broad auditory contextual information (via AC).

\subsection{Limitations}

Our study has limitations including the use of head-fixed mice with restricted behavioral repertoire, the reliance on calcium imaging with its inherent temporal limitations, and the relatively small sample of V2L-injected mice ($n = 3$). The broadband noise stimuli used may not capture the full complexity of natural auditory scenes. Future studies using naturalistic stimuli and freely moving preparations will be important for understanding how these cross-modal projections function during natural behavior.

\section{Conclusion}

We have demonstrated that auditory cortex provides primary visual cortex with robust information about sound source location through direct cortical projections to V1 layer 1. This spatial information is organized non-topographically, without alignment to V1's retinotopic map, and AC boutons represent sounds across both hemispheres of auditory space. Audiovisual modulation of V1 neurons occurs independently of spatial congruency between sound and light stimuli. These findings establish that cross-modal spatial information in cortex follows fundamentally different organizational principles than in subcortical structures, with implications for understanding how the brain integrates multisensory information for perception.

\bibliographystyle{plainnat}
\bibliography{references}

\end{document}
